\chapter{Specifications\label{chap:Specifications}}

\hypertarget{def-globaldeclarationterm}{}
Specifications are lists of global declarations,
which can be one of the following kinds:
\begin{itemize}
  \item Global storage declarations (defined in \chapref{GlobalStorageDeclarations});
  \item Type declarations (defined in \chapref{TypeDeclarations});
  \item Subprogram declarations (defined in \chapref{SubprogramDeclarations}); or
  \item Global pragmas (defined in \chapref{Global Pragma Declarations}).
\end{itemize}

\ExampleDef{Global Declarations}
\listingref{TypecheckDecl} exemplifies various kinds of global declarations ---
types, global storage elements, and subprograms.
All global declarations in \listingref{TypecheckDecl} are well-typed.
\ASLListing{Global declarations}{TypecheckDecl}{\typingtests/TypingRule.TypecheckDecl.asl}

\ChapterOutline
\begin{itemize}
  \item \secref{Formal Relations for Specifications} defines the formal relations used for specifications;
  \item \secref{Formal Relations for Global Declarations} defines the formal relations used for individual
    global declarations;
  \item \secref{Syntax of Specifications} defines the syntax of specifications;
  \item \secref{Abstract Syntax of Specifications} defines the abstract syntax of specifications;
  \item \secref{Typing Specifications} defines the type rules for specifications;
  \item \secref{Overriding Subprograms} defines subprogram overriding and the associated type rules;
  \item \secref{TopologicalOrdering} defines defined-used relations between global declarations,
        which is needed to sort them;
  \item \secref{Dependencies} defines how the global declarations in a specification are sorted;
  \item \secref{Semantics Of Specifications} defines the dynamic semantics of specifications.
\end{itemize}

\section{Formal Relations for Specifications\label{sec:Formal Relations for Specifications}}
\hypertarget{def-specificationterm}{}

\paragraph{Syntax:} Specifications are grammatically derived from $\Nspec$.
\paragraph{Abstract Syntax:} Specifications are represented in abstract syntax by $\spec$.
\paragraph{Typing:} Specifications are annotated by the relation $\typecheckast$, which is defined in
  \TypingRuleRef{TypeCheckAST}.
\paragraph{Semantics:} The dynamic semantics of specifications in given by the relation $\evalspec$,
  which is defined in \SemanticsRuleRef{EvalSpec}.

\section{Formal Relations for Global Declarations\label{sec:Formal Relations for Global Declarations}}
Specifications are lists of global declarations.
The formal relations for individual global declarations are as follows:

\paragraph{Syntax:} Global declarations are grammatically derived from $\Ndecl$.

\paragraph{Abstract Syntax:} Global declarations are represented in abstract syntax by $\decl$.
Global declarations are generated by the relation
\hypertarget{build-decl}{}
\[
  \builddecl(\overname{\parsenode{\Ndecl}}{\vparsednode}) \;\aslrel\; \overname{\KleeneStar{\decl}}{\vastnode}
\]
transforms a parse node $\vparsednode$ into an AST node $\vastnode$.

The case rules for building global declarations are the following:
\begin{itemize}
  \item \ASTRuleRef{GlobalStorageDecl} for global storage declarations;
  \item \ASTRuleRef{TypeDecl} for type declarations;
  \item \ASTRuleRef{SubprogramDecl} for subprogram declarations; and
  \item \ASTRuleRef{GlobalPragma} for global pragmas.
\end{itemize}

\paragraph{Typing:} Global declarations are annotated by the following function.
\RenderRelation{typecheck_decl}

  This function is implemented for each kind of global declarations:
\begin{itemize}
  \item \TypingRuleRef{AnnotateGlobalStorage} for global storage declarations;
  \item \TypingRuleRef{AnnotateTypeDecl} for type declarations; and
  \item \TypingRuleRef{AnnotateSubprogram} for subprogram declarations.
\end{itemize}

We note that, since global pragma declarations are discarded after they are checked
via \TypingRuleRef{CheckGlobalPragma} and do not generate declarations.
Therefore, they have no corresponding rule to implement $\typecheckdecl$.

\paragraph{Semantics:} Only global storage declarations and subprogram declarations are associated
  with a dynamic semantics. Global storage declarations are evaluated via
  $\evalglobals$ and subprogram declarations are evaluated via $\evalcall$.

%%%%%%%%%%%%%%%%%%%%%%%%%%%%%%%%%%%%%%%%%%%%%%%%%%%%%%%%%%%%%%%%%%%%%%%%%%%%%%%%
\section{Syntax of Specifications\label{sec:Syntax of Specifications}}
%%%%%%%%%%%%%%%%%%%%%%%%%%%%%%%%%%%%%%%%%%%%%%%%%%%%%%%%%%%%%%%%%%%%%%%%%%%%%%%%
\begin{flalign*}
\Nspec   \derives\ & \maybeemptylist{\Ndecl} &
\end{flalign*}

%%%%%%%%%%%%%%%%%%%%%%%%%%%%%%%%%%%%%%%%%%%%%%%%%%%%%%%%%%%%%%%%%%%%%%%%%%%%%%%%
\section{Abstract Syntax of Specifications\label{sec:Abstract Syntax of Specifications}}
%%%%%%%%%%%%%%%%%%%%%%%%%%%%%%%%%%%%%%%%%%%%%%%%%%%%%%%%%%%%%%%%%%%%%%%%%%%%%%%%
\RenderType[remove_hypertargets]{spec}


\ASTRuleDef{AST}
\NotImportedToASLSpecYet{
\hypertarget{build-ast}{}
The relation
\[
  \buildast(\overname{\parsenode{\Nspec}}{\vparsednode}) \;\aslrel\; \overname{\spec}{\vastnode}
\]
transforms an $\Nspec$ node $\vparsednode$ into an AST specification node $\vastnode$.
} % END_OF_UNIMPORTED_RELATION

We define this function for subprogram declarations, type declarations, and global storage declarations in the corresponding chapters.

\begin{mathpar}
\inferrule{
    \buildlist[\builddecl](\vdecls) \astarrow \vdeclsone \\
    \concatlist{\vdeclsone} \astarrow \vadecls
}{
    \buildast(\overname{\Nspec(\namednode{\vdecls}{\maybeemptylist{\Ndecl}})}{\vparsednode}) \astarrow \overname{\vadecls}{\vastnode}
}
\end{mathpar}

%%%%%%%%%%%%%%%%%%%%%%%%%%%%%%%%%%%%%%%%%%%%%%%%%%%%%%%%%%%%%%%%%%%%%%%%%%%%%%%%%%%%
\section{Typing Specifications\label{sec:Typing Specifications}}
%%%%%%%%%%%%%%%%%%%%%%%%%%%%%%%%%%%%%%%%%%%%%%%%%%%%%%%%%%%%%%%%%%%%%%%%%%%%%%%%%%%%

The \untypedast{} of an ASL specification consists of a list of global declarations.
Typechecking the untyped AST succeeds if all declarations can be successfully annotated,
which is achieved via \TypingRuleRef{TypeCheckAST}. Otherwise, the result is a
\typingerrorterm{}.

We note that whether typechecking a specification succeeds or fails with a \typingerrorsterm,
does not depend on the order in which the declarations appear in the specification.
That is, if typechecking a specification with one ordering of its declarations leads to
a \typingerrorterm{}, then typechecking that specification with any other ordering of its
declarations also leads to a \typingerrorterm{}, but the \typingerrorsterm{} may not be the same
(since there may be two erroneous declarations and the \typingerrorterm{} returned
depends on which declaration is processed first).

When typechecking declarations, it is important to process them in a certain order,
to avoid false \typingerrorsterm{} resulting from typechecking a declaration that uses an identifier
before the declaration that defines it.
This order relies on the notion of \defusedependencyterm, which we formally define in
\secref{Dependencies}.
\secref{TopologicalOrdering} formally defines how to use the inferred \defusedependenciesterm\ to
order the declarations such that false \typingerrorsterm{} are avoided.


\TypingRuleDef{TypeCheckAST}
\RenderRelation{type_check_ast}

\ExampleDef{Typechecking a List of Global Declarations}
The specification in \listingref{TypeCheckAST} contains numerous global declarations.

\ASLListing{Typechecking a list of global declarations}{TypeCheckAST}{\typingtests/TypingRule.TypeCheckAST.asl}

The typechecker first processes subprogram overrides. Since there are none in this specification,
the list of global subprogram declarations remains the same.
The global declarations are separated into the list of pragmas --- \verb|pragma1| --- and the remaining
global declarations.

The \dependencygraphterm{} for this specification contains a node for each one of
$\vWORDSIZE$, $\MyRecord$, $\vmain$,
$\Color$, $\RED$, $\GREEN$, $\BLUE$,
$\rotatecolor$, $\rotatezero$, $\rotateone$, $\vg$,
and the following edges
\[
\begin{array}{rcl}
  (\Other(\MyRecord)    &,& \Other(\vWORDSIZE))\\
  (\Other(\MyRecord)    &,& \Other(\Color))\\
  (\Other(\RED)       &,& \Other(\Color))\\
  (\Other(\GREEN)       &,& \Other(\Color))\\
  (\Other(\BLUE)       &,& \Other(\Color))\\
  (\Other(\vg)       &,& \Other(\MyRecord))\\
  (\Subprogram(\vmain)  &,& \Other(\vg))\\
  (\Subprogram(\vmain)  &,& \Subprogram(\rotatezero))\\
  (\Subprogram(\rotatecolor)   &,& \Other(\Color))\\
  (\Subprogram(\rotatecolor)   &,& \Other(\RED))\\
  (\Subprogram(\rotatecolor)   &,& \Other(\GREEN))\\
  (\Subprogram(\rotatecolor)   &,& \Other(\BLUE))\\
  (\Subprogram(\rotatezero)    &,& \Other(\Color))\\
  (\Subprogram(\rotatezero)    &,& \Other(\rotatecolor))\\
  (\Subprogram(\rotatezero)    &,& \Other(\rotateone))\\
  (\Subprogram(\rotateone)     &,& \Other(\Color))\\
  (\Subprogram(\rotateone)     &,& \Other(\rotatecolor))\\
  (\Subprogram(\rotateone)     &,& \Other(\rotatezero))\\
\end{array}
\]

Constructing the strongly-connected components for this graph with its edges reversed,
and topologically ordering the strongly-connected components yields the following list of
components:
\[
\begin{array}{l}
  \{ \Other(\vWORDSIZE) \},\\
  \{ \Other(\MyRecord) \},\\
  \{ \Other(\vg) \},\\
  \{ \Other(\Color) \},\\
  \{ \Other(\RED) \},\\
  \{ \Other(\GREEN) \},\\
  \{ \Other(\BLUE) \},\\
  \{ \Subprogram(\rotatecolor) \},\\
  \{ \Subprogram(\rotatezero), \Subprogram(\rotateone) \},\\
  \{ \Subprogram(\vmain) \}
\end{array}
\]
The only non-trivial component (that is, a component with more than one global declaration)
is $\{ \Subprogram(\rotatezero), \Subprogram(\rotateone) \}$, since these are the only
mutually-recursive subprograms in this specification.

The typechecker annotates each strongly-connected component listed above,
and finally checks that the pragma \verb|pragma pragma1;| is well-typed, which it is.

\RenderProseAndFormally{type_check_ast}
\CodeSubsection{\TypeCheckASTBegin}{\TypeCheckASTEnd}{../Typing.ml}

\TypingRuleDef{AnnotateDeclComps}
\RenderRelation{annotate_decl_comps}

We note that a strongly-connected component containing just a single declaration may contain
any kind of global declaration ---
a type declaration, a global storage declaration, or a subprogram declaration ---
whereas a strongly-connected component containing multiple declarations must be checked
to contain only subprograms. This is because the only type of mutually-recursive declarations
allowed in ASL are between subprograms. The rules below handle these cases separately (\textsc{single}
for single declarations and \textsc{mutually\_recursive} for more than one declaration).

\ExampleDef{Annotating Strongly-connected Components of Declarations}
In \ExampleRef{Typechecking a List of Global Declarations},
each declaration is in its own strongly-connected component,
except for the component
\[
\{ \Subprogram(\rotatezero), \Subprogram(\rotateone) \} \enspace.
\]

\RenderProseAndFormally{annotate_decl_comps}

\TypingRuleDef{TypeCheckMutuallyRec}
\RenderRelation{type_check_mutually_rec}

\ExampleDef{Checking Mutually-recursive Declarations}
In \ExampleRef{Typechecking a List of Global Declarations},
the strongly-connected component
$\{ \Subprogram(\rotatezero), \Subprogram(\rotateone) \}$
contains only subprogram declarations.

In the specification in \listingref{TypeCheckMutuallyRec-bad},
the global storage element \verb|g| is defined in terms of the
subprogram \verb|foo|, which is defined in terms of \verb|g|.
%
Technically, the \dependencygraphterm{} contains the strongly-connected
component $\{\Other(\vg), \Subprogram(\foo)\}$.
%
Since mutual recursion is allowed only between subprograms,
this specification is ill-typed.

\ASLListing{Illegal recursion among global declarations}{TypeCheckMutuallyRec-bad}{\typingtests/TypingRule.TypeCheckMutuallyRec.bad.asl}

The specification in \listingref{TypeCheckMutuallyRec-bad2} is ill-typed,
since it contains recursive type declarations.
\ASLListing{Illegal recursive type declarations}{TypeCheckMutuallyRec-bad2}{\typingtests/TypingRule.TypeCheckMutuallyRec.bad2.asl}

\RenderProseAndFormally{type_check_mutually_rec}
\CodeSubsection{\TypeCheckMutuallyRecBegin}{\TypeCheckMutuallyRecEnd}{../Typing.ml}

\TypingRuleDef{FuncsOfDecls}
\RenderRelation{funcs_of_decls}
See \ExampleRef{Checking Mutually-recursive Declarations}.
\RenderProseAndFormally{funcs_of_decls}

\TypingRuleDef{DeclareSubprograms}
\RenderRelation{declare_subprograms}

\ExampleDef{Updating Static Environments for Subprogram Declarations}
Applying $\declaresubprograms$ for the subprogram declarations in \listingref{DeclareSubprograms}
adds the following bindings to the $\overloadedsubprograms$ map
\[
\begin{array}{rcl}
  \flipbits &\mapsto& \{\flipbits\} \\
  \addten    &\mapsto& \{\addten, \addtenone\} \\
  \factorial  &\mapsto& \{\factorial\} \\
\end{array}
\]
and updates the $\subprograms$ map by binding
strings representing unique subprogram names to the $\func$ nodes corresponding to
the subprogram signatures:
\begin{center}
\begin{tabular}{ll}
\textbf{string}     & \textbf{signature}\\
\texttt{flip\_bits} & \verb|func flip_bits{N}(bv: bits(N)) => bits(N)|\\
\texttt{add\_10}    & \verb|func add_10(x: real) => real|\\
\texttt{add\_10-1}  & \verb|func add_10(x: integer) => integer|\\
\texttt{factorial}  & \verb|func factorial(x: integer) => integer recurselimit 100|
\end{tabular}
\end{center}

Using the name \texttt{add\_10-1} for \verb|func add_10(x: integer) => integer|
and \texttt{add\_10} for \verb|func add_10(x: real) => real| (rather than
the other way around) is due to the arbitrary choice of a topological ordering
of the subprogram declarations.

\ASLListing{Updating static environments for subprogram declarations}{DeclareSubprograms}{\typingtests/TypingRule.DeclareSubprograms.asl}

\RenderProseAndFormally{declare_subprograms}

\TypingRuleDef{AddSubprogramDecls}
\RenderRelation{add_subprogram_decls}

\ExampleDef{Adding Subprogram Declarations}
The specification in \listingref{AddSubprogramDecls}
contains two subprogram declarations,
which are added to the \staticenvironmentterm{}.
The subprogram with the \integertypeterm{} argument is named \verb|increment|
and the subprogram with the \realtypeterm{} argument is named \verb|increment-1|.
\ASLListing{Adding subprogram declarations}{AddSubprogramDecls}{\typingtests/TypingRule.AddSubprogramDecls.asl}

\RenderProseAndFormally{add_subprogram_decls}

\section{Overriding Subprograms\label{sec:Overriding Subprograms}}
This section defines how to override subprograms in a specification.
In particular, a subprogram marked by \texttt{impdef} can be overridden by a subprogram marked by \\
\texttt{implementation}.

\RequirementDef{OverridingMatch}
An overriding \Proseimplementationsubprogram{} subprogram \\
must match the overridden \Proseimpdefsubprogram{} subprogram in all of the following:
\begin{itemize}
  \item name of subprogram;
  \item qualifier;
  \item number, names, types, and order of arguments;
  \item number, names, types, and order of parameters;
  \item return type.
\end{itemize}
In \ExampleRef{Overriding Subprograms}, the \verb|Foo| subprograms marked with \texttt{impdef} and the one marked with \texttt{implementation} match in all of the above.

\ExampleDef{Overriding Subprograms}
\listingref{overriding} shows a specification which defines \Proseimpdefsubprograms{} \verb|Foo| and \verb|Bar|.
The definition of \verb|Foo| is overridden by a corresponding \Proseimplementationsubprogram{}.
\ASLListing{Overriding subprograms}{overriding}{\definitiontests/Overriding.asl}

\noindent
Subprograms marked with \texttt{implementation} are subject to the following restrictions.

\RequirementDef{NonClashingImplementations}
No two \Proseimplementationsubprograms{} can \linebreak[3]
match each other according to \RequirementRef{OverridingMatch}.
In \ExampleRef{Invalid Overriding}, the \verb|Foo| subprograms marked \texttt{implementation}
match and are therefore invalid.

\RequirementDef{SingleOverrideCandidate}
Each \Proseimplementationsubprogram{} must override exactly one \Proseimpdefsubprogram.
In \ExampleRef{Invalid Overriding}, the \Proseimplementationsubprogram{} \verb|Bar| is
invalid as it has no corresponding \Proseimpdefsubprogram{}.

\ExampleDef{Invalid Overriding}
\listingref{overriding-bad} shows a specification which attempts to define clashing
\Proseimplementationsubprograms{} named \verb|Foo|, and define an
\Proseimplementationsubprogram{} \verb|Bar| without a corresponding \Proseimpdefsubprogram{}.
\ASLListing{Invalid overriding}{overriding-bad}{\definitiontests/OverridingBad.asl}

\noindent
Implementations may wish to emit user-configurable warnings to check either:
\begin{itemize}
  \item All \Proseimpdefsubprograms{} have been overridden by corresponding \Proseimplementationsubprograms{}.
  \item There are no \Proseimplementationsubprograms{}, so no \Proseimpdefsubprograms{} are overridden.
\end{itemize}

\TypingRuleDef{OverrideSubprograms}
\RenderRelation{override_subprograms}
See \ExampleRef{Overriding Subprograms} for an example of valid overriding.

\RenderProseAndFormally{override_subprograms}

\TypingRuleDef{OverrideDeclsSort}
\RenderRelation{override_decls_sort}
See \ExampleRef{Overriding Subprograms} for an example of valid overriding.

\RenderProseAndFormally{override_decls_sort}

\TypingRuleDef{CheckImplementationsUnique}
\RenderRelation{check_implementations_unique}
See \ExampleRef{Invalid Overriding} for an example of non-unique \Proseimplementationsubprograms{}.

\RenderProseAndFormally{check_implementations_unique}

\TypingRuleDef{SignaturesMatch}
\RenderRelation{signatures_match}
See \ExampleRef{Overriding Subprograms} for an example of matching signatures.

\RenderProseAndFormally{signatures_match}

\TypingRuleDef{ProcessOverrides}
\RenderRelation{process_overrides}
See \ExampleRef{Overriding Subprograms} for an example of valid replacement of an \Proseimplementationsubprogram{}.

\RenderProseAndFormally{process_overrides}

\TypingRuleDef{RenameSubprograms}
\RenderRelation{rename_subprograms}

\ExampleDef{Typechecking overridden subprograms}
The specification in \listingref{RenameSubprograms} contains an ill-typed \Proseimpdefsubprogram{} that is overridden.
This is given a fresh name by $\renamesubprograms$, so that it is still typechecked but never referred to in the resulting specification.
In particular, though the ill-typed subprogram is overridden and will not be used in the annotated specification, it still results in a
\typingerrorterm.
\ASLListing{Typechecking an overridden subprogram}{RenameSubprograms}{\typingtests/TypingRule.RenameSubprograms.asl}

\RenderProseAndFormally{rename_subprograms}

\section{Defined-Used Dependencies for Global Declarations\label{sec:TopologicalOrdering}}
This section defines how to construct a graph of \defusedependenciesterm\ between the identifiers associated
with global declarations.
This is achieved by associating, for each declaration $d$, the set identifiers \emph{used} by $d$,
which is formally defined by $\usedecl$, and the identifier \emph{defined} by $d$,
which is formally defined by $\defdecl$.
%
The set of \defusedependenciesterm\ associated with a declaration $\vd$ is given by
$\decldependencies(\vd)$ (see \TypingRuleRef{DeclDependencies}).

\RequirementDef{GlobalNamespace}
The global namespace effectively consists of two independent namespaces: one for subprogram names, and the other for all other identifiers.

\ExampleDef{Global Namespace}
\listingref{globalnamespace} shows a specification which defines a subprogram \verb|X|
that does not interfere with the variable declaration also named \verb|X|.
\ASLListing{Global namespace}{globalnamespace}{\definitiontests/GlobalNamespace.asl}

We distinguish between subprogram identifiers, and all other identifiers (storage elements and \namedtypes{}).
We use the following type to track \defusedependenciesterm{} between subprogram identifiers and other identifiers:
\RenderType{def_use_name}


\TypingRuleDef{BuildDependencies}
\RenderRelation{build_dependencies}

\ExampleDef{The Dependency Graph for a List of Global Declarations}
The \dependencygraphterm{} generated for the specification in \listingref{DeclDependencies}
is as follows:
\[
\left(
\begin{array}{c}
  \{ \Other(\vg), \Other(\MyRecord), \Other(\Color), \Other(\vWORDSIZE), \Subprogram(\vmain) \},\\
  \left\{\begin{array}{rcl}
    ( \Other(\vg) &,&          \Other(\MyRecord) ),\\
    ( \Other(\MyRecord) &,&    \Other(\vWORDSIZE) ),\\
    ( \Other(\Color) &,&       \Other(\RED) ),\\
    ( \Other(\Color) &,&       \Other(\GREEN) ),\\
    ( \Other(\Color) &,&       \Other(\BLUE) ),\\
    ( \Subprogram(\vmain) &,&  \Other(\vg) )\\
  \end{array}\right\}
\end{array}
\right)
\]

\RenderProseAndFormally{build_dependencies}

\TypingRuleDef{DeclDependencies}
\RenderRelation{decl_dependencies}

\ExampleDef{The Dependencies of Global Declarations}
The specification in \listingref{DeclDependencies}
shows the dependencies generated for each global declarations
in comments appearing to the right of them or just above them.
A dependency $(d_1,d_2)$ is depicted as \texttt{$d_1$ -> $d_2$}.
\ASLListing{Dependencies of global declarations}{DeclDependencies}{\typingtests/TypingRule.DeclDependencies.asl}

\RenderProseAndFormally{decl_dependencies}

\TypingRuleDef{DefDecl}
\RenderRelation{def_decl}

\ExampleDef{The Identifiers Defined by Global Declarations}
The specification given in \listingref{DefDecl},
shows examples of global declarations and the identifiers they define,
which appear in comments to their right.
\ASLListing{Identifiers defined by global declarations}{DefDecl}{\typingtests/TypingRule.DefDecl.asl}

\RenderProseAndFormally{def_decl}

\TypingRuleDef{DefEnumLabels}
\RenderRelation{def_enum_labels}

\ExampleDef{Identifiers Defined by Declaring an Enumeration Type}
The set of identifiers defined by the enumeration declaration
in \listingref{DefEnumLabels} is \\
$\{\Other(\RED), \Other(\GREEN), \Other(\BLUE)\}$.
\ASLListing{Identifiers defined by declaring an enumeration type}{DefEnumLabels}{\typingtests/TypingRule.DefEnumLabels.asl}

\RenderProseAndFormally{def_enum_labels}

\TypingRuleDef{UseDecl}
\RenderRelation{use_decl}

\ExampleDef{Identifiers Used by Global Declarations}
The specification in \listingref{UseDecl} shows the set of identifiers
used by each global declaration via comments above it.
\ASLListing{Identifiers used by global declarations}{UseDecl}{\typingtests/TypingRule.UseDecl.asl}

\RenderProseAndFormally{use_decl}

\TypingRuleDef{UseTy}
\RenderRelation{use_ty}

\ExampleDef{The Identifiers Used by a Type}
The specification in \listingref{UseTy} shows type annotations and the identifiers
used by them appearing as comments to their right.
\ASLListing{The identifiers used by a type}{UseTy}{\typingtests/TypingRule.UseTy.asl}

\RenderProseAndFormally{use_ty}

\TypingRuleDef{UseSubtypes}
\RenderRelation{use_subtypes}

\ExampleDef{The Identifiers Used by a Subtyping Declaration}
In \listingref{UseDecl}, applying $\usesubtypes$ at
the type declaration of \verb|MyRecord| yields \\
$\{\ \Other(\RecordBase), \Other(\HalfWordBits)\ \}$.

\RenderProseAndFormally{use_subtypes}

\TypingRuleDef{UseExpr}
\RenderRelation{use_expr}

\ExampleDef{The Identifiers Used by Expressions}
The specification in \listingref{UseExpr} shows the identifiers used
by expressions on the right-hand-side of assignments in comments appearing to their right.
\ASLListing{Identifiers used by expressions}{UseExpr}{\typingtests/TypingRule.UseExpr.asl}

\RenderProseAndFormally{use_expr}

\TypingRuleDef{UseLexpr}
\RenderRelation{use_lexpr}

\ExampleDef{The Identifiers Used by Assignable Expressions}
The specification in \listingref{UseLexpr} shows the identifiers used
by the assignable expressions on the left-hand-side of assignments in comments appearing to their right.
\ASLListing{Identifiers used by assignable expressions}{UseLexpr}{\typingtests/TypingRule.UseLexpr.asl}

\RenderProseAndFormally{use_lexpr}

\TypingRuleDef{UsePattern}
\RenderRelation{use_pattern}

\ExampleDef{The Identifiers Used by a Pattern}
The specification in \listingref{UsePattern} shows examples of patterns
and the identifiers used by them, appearing in comments to their right.
\ASLListing{Identifiers used by a pattern}{UsePattern}{\typingtests/TypingRule.UsePattern.asl}

\RenderProseAndFormally{use_pattern}

\TypingRuleDef{UseSlice}
\RenderRelation{use_slice}

\ExampleDef{The Identifiers Used by a Slice}
The specification in \listingref{UseSlice} shows slicing expressions
and the identifiers they use, appearing in comments to their right.
\ASLListing{Identifiers used by a slice}{UseSlice}{\typingtests/TypingRule.UseSlice.asl}

\RenderProseAndFormally{use_slice}

\TypingRuleDef{UseBitfield}
\RenderRelation{use_bitfield}

\ExampleDef{The Identifiers Used by a Bitfield Declaration}
The specification in \listingref{UseBitfield} shows examples
of bitfield declarations and the identifiers they use,
appearing as comments to their right.
\ASLListing{Identifiers used by a bitfield declaration}{UseBitfield}{\typingtests/TypingRule.UseBitfield.asl}

\RenderProseAndFormally{use_bitfield}

\TypingRuleDef{UseConstraint}
\RenderRelation{use_constraint}

\ExampleDef{The Identifiers Used by Constraints}
The specification in \listingref{UseConstraint}
shows constraints in type annotations, and the identifiers they use,
appearing in comments to their right.
\ASLListing{Identifiers used by constraints}{UseConstraint}{\typingtests/TypingRule.UseConstraint.asl}

\RenderProseAndFormally{use_constraint}

\TypingRuleDef{UseLDI}
\RenderRelation{use_ldi}

\ExampleDef{The identifiers used by local declaration items}
The specification in \listingref{UseLDI} shows examples of local declaration items
and the identifiers used by them, appearing in the comments to their right.
\ASLListing{Identifiers used by local declaration items}{UseLDI}{\typingtests/TypingRule.UseLDI.asl}

\RenderProseAndFormally{use_ldi}

\TypingRuleDef{UseStmt}
\RenderRelation{use_stmt}

\ExampleDef{The Identifiers Used by Statements}
The specification in \listingref{UseStmt} shows examples of statements
and the identifiers used by them, appearing in comments to their right.
\ASLListing{Identifiers used by statements}{UseStmt}{\typingtests/TypingRule.UseStmt.asl}

\RenderProseAndFormally{use_stmt}

\TypingRuleDef{UseCatcher}
\RenderRelation{use_catcher}
\ExampleDef{The Identifiers Used by a Catcher}
The specification in \listingref{UseStmt} shows an example
of a \catcherterm{} (towards the end of \verb|main|)
and the identifiers used by it, appearing in comments to its right.

\RenderProseAndFormally{use_catcher}

\section{Ordering Global Declarations via Defined-Used Dependencies\label{sec:Dependencies}}

\begin{definition}[Strongly Connected Components]
\hypertarget{relation-scc}{}
Given a graph $G=(V, E)$, a \\ subset of its nodes $C \subseteq V$ is called
a \emph{strongly connected component} of $G$ if
every pair of nodes $u,v \in C$ reachable from one another.

The \emph{strongly connected components} of a graph $(V, E)$ uniquely partitions its set of
nodes $V$ into a set of strongly connected components:
\[
\sccop[H]{V}{E} \triangleq \{ C \subseteq V \;|\; \forall u,v\in C.\ (u,v), (v,u) \in \graphtransitivereflexive{E} \} \enspace.
\]
\end{definition}

\hypertarget{relation-topologicalorderingcomps}{}
\begin{definition}[Topological Ordering of Components]
For a non-empty graph \\
$G=(V, E)$ and its strongly connected components $\comps \triangleq \sccop[H]{V}{E}$,
a listing of $\comps$ ---
$C_{1..k}$ --- is a \emph{topological ordering of components},
denoted \\
$C_{1..k} \in \topologicalorderingcomps[H]{\comps}{E}$, if the following condition holds:
\[
  \forall 1 \leq i \leq j \leq k.\ \exists c_i\in C_i.\ c_j\in C_j.\ (c_i,c_j) \in \graphtransitivereflexive{E} \;\;\implies\;\;
  i \leq j \enspace .
\]
\end{definition}

\section{Semantics of Specifications\label{sec:Semantics Of Specifications}}
The semantics of specifications is defined via the relation $\evalspec$, which is defined next.

\SemanticsRuleDef{EvalSpec}
\RenderRelation{eval_spec}

\ExampleRef{Returning a Value from the Entry Point} shows a specification
whose evaluation results in returning the value $\nvint(0)$.

\ExampleRef{An Uncaught Exception} shows a specification whose evaluation
results in a \dynamicerrorterm.

\RenderProseAndFormally{eval_spec}
\CodeSubsection{\EvalSpecBegin}{\EvalSpecEnd}{../Interpreter.ml}

\SemanticsRuleDef{BuildGlobalEnv}
\RenderRelation{build_genv}

It is assumed that $\typedspec$ lists the declarations in reverse order with respect
to the \defusedependencyterm\ order
(see \TypingRuleRef{TypeCheckAST}).

See \ExampleRef{Evaluating a List of Global Declarations}.

\RenderProseAndFormally{build_genv}
\CodeSubsection{\EvalBuildGlobalEnvBegin}{\EvalBuildGlobalEnvEnd}{../Interpreter.ml}
